\section{Running it}
\label{sec:estimation}

\begin{algorithm}[t]
\caption{The forecast, as a trainer would run it}
\label{alg:forecast}
\small
\begin{algorithmic}[1]
  \State train as usual, storing $\theta_t$ and $\eta_t$ for each update
  \State $b \gets \nabla M(\theta_T)$ \Comment{one rollout and backward pass on the eval set}
  \State $v \gets 0$
  \For{$t = T-1$ \textbf{down to} $0$}
    \State draw a batch at $\theta_t$; let $\hat{g}_{t,i}$ be its per-prompt gradients
    \State $v \gets v + (\eta_t^2/P)\,\widehat{\Var}_i\big(b^{\!\top}\hat{g}_{t,i}\big)$
      \Comment{one scalar}
    \If{transport is needed}
      \State $\delta \gets \epsilon\, b/\|b\|$
      \State $b \gets b + \dfrac{\eta_t}{2}\cdot\dfrac{\|b\|}{\epsilon}
        \big(\gbar(\theta_t + \delta) - \gbar(\theta_t - \delta)\big)$
        \Comment{common random numbers}
    \EndIf
  \EndFor
  \State \Return $\sqrt{v}$
\end{algorithmic}
\end{algorithm}

\paragraph{What it costs, plainly.} As written, the backward pass makes three passes per stored
update: one batch for the injected term and two for the Hessian-vector product. That is about three
times the compute of the training run, not a rounding error, and we would rather state it than let
a reader discover it.

It is also usually unnecessary. Ablating the transport, holding $b$ at $\nabla M(\theta_T)$ instead
of carrying it back, changes the forecast by a median of $2\%$ and the median error not at all over
runs of ten and twenty-five updates. The reason is in the spectrum of $A_t$: most of it sits within
$10^{-4}$ of one and $\nabla M$ points into that part, so $\Phi_{T,t}$ is close to the identity
along the one direction the forecast needs. What we recommend is the cheap estimator,
\begin{equation}
  \label{eq:cheap}
  \Var_{\text{seed}}\big[M(\theta_T)\big] \;\approx\;
  \sum_{t} \frac{\eta_t^2}{P}\,\widehat{\Var}_i\!\big(\nabla M(\theta_T)^{\!\top}\hat{g}_{t,i}\big),
\end{equation}
one extra pass per update and no Hessian-vector products, with the full recursion held in reserve
for runs long enough, or curvature strong enough, that a per-update discrepancy of a percent
compounds. Knowing which regime one is in is what the recursion is for.

\paragraph{A closed form with no run at all.} Take \eqref{eq:cheap} with $\Phi \approx I$ and a
constant step, replace the projection $\Var_i(\nabla M^{\!\top}\hat g_{t,i})$ by the trace
\eqref{eq:two-level}, and the run drops out entirely:
\begin{equation}
  \label{eq:frozen}
  \sigma \;\propto\; \eta \sqrt{T/P}\,\sqrt{\tau_{\mathrm{b}} + \tau_{\mathrm{w}}/(G-1)}.
\end{equation}
We call this the \emph{frozen} forecast because it needs only the run length, the batch shape and
the two noise traces, and is therefore available before any training happens.

It is written as a proportionality on purpose. The constant it is missing carries
$\|\nabla M\|$, which is exactly the part that requires a trained model, and on our grid
\eqref{eq:frozen} overstates the measured deviation by a factor of about
$\numSizeFrozenScale$. That factor is not quite constant: it varies by
$\numSizeFrozenScaleSpread\times$ across the grid, leaving a relative error of
$\numSizeFrozenRelative\times$ at the median once one common scale is removed. So
\eqref{eq:frozen} is useless for reporting an error bar and, by the argument of
Section~\ref{sec:sizing}, entirely serviceable for sizing one, because a stopping threshold absorbs
a constant and nothing absorbs the part that varies.

Both traces come from the microbatch gradients a trainer already accumulates: with blocks $a$ holding
disjoint prompts and sub-groups $b \in \{1,2\}$ independent given a prompt,
\begin{equation}
  \label{eq:estimators}
  \begin{aligned}
  \tr\Sb &= \tfrac{P}{K}\big( \E[ g_{a1}^{\!\top} g_{a2} ] - \E[ g_{a1}^{\!\top} g_{a'2} ]_{a \neq a'} \big),\\
  \tr\Sw &= \tfrac{P}{2K}\, \E\big[ \| g_{a1} - g_{a2} \|^2 \big].
  \end{aligned}
\end{equation}
The cost is $2K$ gradient buffers rather than one, no extra responses and no extra backward work.
Estimating the total and subtracting the within-prompt part instead differences two large
correlated quantities and is badly conditioned: on homogeneous pools it returns a negative
$\tau_{\mathrm{b}}$. Against exactly known values over pools whose noise ratio spans a factor of
$26$, \eqref{eq:estimators} recovers the optimal group size to a median error of $1.3\%$, while the
subtractive form reports $6.1$ where the truth is $9.4$.

Section~\ref{sec:sizing} finds that \eqref{eq:frozen} does most of the work, which matters because
it costs almost nothing.
